\section{Results}
\label{sec:results}

We organize results around four research questions (RQ1--RQ4) using a downstream next-event prediction task. Unless otherwise stated, models are evaluated on the held-out real test split (\texttt{real\_test.npz}), and we report macro-F1 as the primary metric, with accuracy and Top-5 accuracy as secondary metrics.

\subsection{RQ1: Can synthetic traces replace real data for downstream training?}
\label{sec:rq1}

RQ1 evaluates the utility of synthetic traces for training downstream models by comparing four training regimes: real-only (baseline), synthetic-only (raw), synthetic-only (repaired), and a 50/50 mixture of real+synthetic (augmentation). All models are tested on the same held-out real test set.

\begin{table}[t]
\centering
\caption{RQ1 (Data utility): next-event prediction performance when training on real vs. synthetic data (dummy values shown).}
\label{tab:rq1}
\begin{tabular}{lccc}
\toprule
\textbf{Training data} & \textbf{Macro-F1} & \textbf{Accuracy} & \textbf{Top-5 Acc.} \\
\midrule
Real (baseline)                & 0.920 & 0.945 & 0.985 \\
Synthetic (raw, $L{=}1024$)    & 0.710 & 0.790 & 0.910 \\
Synthetic (repaired, $L{=}1024$) & 0.840 & 0.885 & 0.955 \\
Real + Synthetic (50/50, $L{=}1024$) & 0.930 & 0.950 & 0.988 \\
\bottomrule
\end{tabular}
\end{table}

\noindent\textbf{Takeaway (template):} Repaired synthetic traces retain $\sim$XX--YY\% of the real-data baseline macro-F1, and mixing repaired synthetic with real data provides modest gains over the real-only baseline.

\subsection{RQ2: Does constraint-guided repair improve downstream performance?}
\label{sec:rq2}

RQ2 isolates the impact of repair by comparing models trained on raw synthetic traces versus repaired synthetic traces, across multiple context lengths. We report the absolute and relative improvements in macro-F1.

\begin{table}[t]
\centering
\caption{RQ2 (Repair effectiveness): improvement from repair across context lengths (dummy values shown).}
\label{tab:rq2}
\begin{tabular}{lcccc}
\toprule
\textbf{Context $L$} & \textbf{Raw F1} & \textbf{Repaired F1} & \textbf{$\Delta$F1} & \textbf{Rel. gain} \\
\midrule
256  & 0.660 & 0.790 & +0.130 & +19.7\% \\
1024 & 0.710 & 0.840 & +0.130 & +18.3\% \\
4096 & 0.740 & 0.860 & +0.120 & +16.2\% \\
\bottomrule
\end{tabular}
\end{table}

\noindent\textbf{Takeaway (template):} Constraint-guided repair consistently improves utility, yielding a macro-F1 gain of +0.12 to +0.13 across context lengths, indicating that correcting structural invalidities translates into measurable downstream benefits.

\subsection{RQ3: Does context length affect synthetic trace utility?}
\label{sec:rq3}

RQ3 evaluates whether longer context windows used during diffusion training lead to higher-quality synthetic traces for downstream learning. To control for validity effects, we compare repaired synthetic traces generated by models trained with $L \in \{256, 1024, 4096\}$.

\begin{table}[t]
\centering
\caption{RQ3 (Context length): next-event prediction when training on repaired synthetic traces from different diffusion context lengths (dummy values shown).}
\label{tab:rq3}
\begin{tabular}{lccc}
\toprule
\textbf{Context $L$} & \textbf{Macro-F1} & \textbf{Accuracy} & \textbf{Top-5 Acc.} \\
\midrule
256  & 0.790 & 0.845 & 0.945 \\
1024 & 0.840 & 0.885 & 0.955 \\
4096 & 0.860 & 0.900 & 0.962 \\
\bottomrule
\end{tabular}
\end{table}

\noindent\textbf{Takeaway (template):} Increasing diffusion context length improves downstream utility, with $L{=}4096$ outperforming $L{=}256$ by $\Delta$F1 $\approx$ 0.07 (dummy), suggesting that longer-range context helps capture dependencies beneficial for prediction.

\subsection{RQ4: Which input channels are most critical for downstream utility?}
\label{sec:rq4}

RQ4 measures feature importance for the downstream predictor by comparing an event-only predictor to a full-feature predictor trained on real traces. This isolates how much additional system metadata (e.g., \texttt{dt}, \texttt{cpu}, \texttt{tid}, \texttt{comm}, \texttt{ret}) contributes to predictability and task utility.

\begin{table}[t]
\centering
\caption{RQ4 (Feature ablation): event-only vs full-feature next-event prediction on real traces (dummy values shown).}
\label{tab:rq4}
\begin{tabular}{lcc}
\toprule
\textbf{Predictor input} & \textbf{Macro-F1} & \textbf{Accuracy} \\
\midrule
Event-only              & 0.770 & 0.820 \\
Full features           & 0.920 & 0.945 \\
\bottomrule
\end{tabular}
\end{table}

\noindent\textbf{Takeaway (template):} Full-feature inputs substantially outperform event-only inputs, confirming that timing and scheduling metadata provide critical predictive signal and therefore represent an important target for faithful synthetic generation.